
\chapter*{Introduction Générale}
\markboth{Introduction Générale}{} %pour afficher l'entete
\addcontentsline{toc}{chapter}{Introduction Générale}

Le harcèlement scolaire touche des centaines de milliers d'élèves chaque année.
Avec la généralisation d'Internet et des réseaux sociaux, il prend désormais une dimension numérique : le cyberharcèlement.
Les victimes reçoivent des insultes, des menaces et des humiliations en ligne, souvent de manière anonyme,
sans qu'aucun mécanisme ne soit disponible pour les détecter ou les accompagner en dehors des heures ouvrées.

\bigskip

Face à ce constat, la société Smart-Conseil a développé \textbf{Netethic}, une plateforme numérique dédiée
à la détection et à la prévention du harcèlement en milieu scolaire.
Notre projet de fin d'études s'inscrit dans le développement de cette plateforme.
Il porte sur la conception et la réalisation de quatre modules basés sur l'intelligence artificielle,
permettant d'automatiser la détection des cas préoccupants et d'offrir un support aux utilisateurs
disponible à toute heure.

\bigskip

Ce travail a été réalisé en binôme au sein de Smart-Conseil, en suivant la méthode Scrum.
Il réunit deux axes complémentaires : un ensemble de modules conversationnels (assistant Noa, assistant technique, assistant émotionnel)
et un système d'analyse automatique des textes (classificateur de harcèlement).

\bigskip

Ce rapport est organisé en quatre chapitres :

\begin{itemize}
  \item \textbf{Chapitre 1 — Cadre général et Spécification des besoins :}
        Présente l'entreprise Smart-Conseil, analyse les solutions existantes et leurs limites,
        décrit la solution proposée et la méthodologie adoptée, puis définit les besoins
        fonctionnels et non fonctionnels de chaque module ainsi que les acteurs du système.
  \item \textbf{Chapitre 2 — État de l'art et choix technologiques :}
        Compare les approches existantes pour chaque domaine technique couvert
        et justifie les technologies retenues.
  \item \textbf{Chapitre 3 — Conception et Mise en place :}
        Présente l'architecture commune du système, la conception spécifique de chaque module
        illustrée par des diagrammes UML, la construction des bases de connaissances RAG
        et le cycle complet d'apprentissage du classificateur de harcèlement.
  \item \textbf{Chapitre 4 — Réalisation et tests :}
        Présente les interfaces réalisées, les résultats des tests automatisés, l'évaluation
        de la qualité des réponses, les mesures de sécurité selon OWASP, les optimisations
        de performance et le déploiement du système.
\end{itemize}
